% GENERATED FILE. Edit canonical lessons, then run npm run generate:books.
% canonical-source-hash: fnv1a64:a622c69d83b99a9f
% canonical-lessons: JA-C199-shuppatsu, JA-C199-touchaku, JA-C199-tokkyuu, JA-C199-tomaru, JA-C199-kyuukou

\chapter{Departures, Arrivals and Expresses}
\label{ch:ja-departures-arrivals-and-expresses}
\hlchaptermodality{199}

\begin{chapteropening}
\textbf{By the end of this chapter:} \emph{I can read when trains leave, arrive and stop.}

Complete the last lesson of chapter 199: I can read when trains leave, arrive and stop.
\end{chapteropening}

\section[\texorpdfstring{\ja{しゅっぱつ} (shuppatsu)}{しゅっぱつ (shuppatsu)}]{\ja{しゅっぱつ} (shuppatsu) — departure}

\label{lesson:JA-C199-shuppatsu}



\begin{quote}
Before the new one: say the Japanese for a subway, then the Japanese for to get to.
\end{quote}

\subsection*{You'll want to know: \ja{しゅっぱつ}}
\textbf{\ja{しゅっぱつ}} — \emph{shuppatsu} — \textquotedblleft{}departure\textquotedblright{}.

\begin{quote}
\textbf{\ja{しゅっぱつは} \ja{ごごです。}} — \emph{shuppatsu wa gogo desu} — \textquotedblleft{}Departure is in the afternoon.\textquotedblright{}
\end{quote}

\begin{cousinweb}
\textbf{\ja{しゅっぱつ}} (\emph{shuppatsu}) is Sino-Japanese (kango): two Chinese characters read with their on'yomi, readings borrowed from Chinese. They are \emph{shutsu}, \textquotedblleft{}to go out\textquotedblright{} (Middle Chinese \emph{tsyhwit}) and \emph{hatsu}, \textquotedblleft{}to set off\textquotedblright{} (\emph{pjot}) (KANJIDIC2; Middle Chinese in Baxter's transcription). The two join as \emph{shup-patsu}.
\end{cousinweb}

\begin{grammarlens}[title={What you've built}]
One more word for everyday life at A2.
\end{grammarlens}

\subsection*{Your turn}
\begin{itemize}
\raggedright
  \item \emph{Say it:} \emph{shuppatsu}
  \item \emph{Say it:} \emph{shuppatsu}, once more
  \item \emph{Say it:} \emph{shuppatsu wa gogo desu} — Departure is in the afternoon.
  \item \emph{Recall:} say the Japanese for a subway, then the Japanese for to get to, then say \emph{shuppatsu} again
\end{itemize}

\begin{tcolorbox}[breakable,colback=teal!4,colframe=teal!35!black,title={Before you move on}]
What does \ja{しゅっぱつ} mean? (\textquotedblleft{}Departure\textquotedblright{}.) Say it once more.
\end{tcolorbox}
\section[\texorpdfstring{\ja{とうちゃく} (tōchaku)}{とうちゃく (tōchaku)}]{\ja{とうちゃく} (tōchaku) — arrival}

\label{lesson:JA-C199-touchaku}



\begin{quote}
Before the new one: say the Japanese for to get to, then the Japanese for departure.
\end{quote}

\subsection*{You'll want to know: \ja{とうちゃく}}
\textbf{\ja{とうちゃく}} — \emph{tōchaku} — \textquotedblleft{}arrival\textquotedblright{}.

\begin{quote}
\textbf{\ja{とうちゃくは} \ja{ゆうがたです。}} — \emph{tōchaku wa yūgata desu} — \textquotedblleft{}Arrival is in the evening.\textquotedblright{}
\end{quote}

\begin{cousinweb}
\textbf{\ja{とうちゃく}} (\emph{tōchaku}) is Sino-Japanese (kango): two Chinese characters read with their on'yomi, readings borrowed from Chinese. They are \emph{tō}, \textquotedblleft{}to reach\textquotedblright{} (Middle Chinese \emph{tawH}) and \emph{chaku}, \textquotedblleft{}to arrive\textquotedblright{} (\emph{drjak}) (KANJIDIC2; Middle Chinese in Baxter's transcription).
\end{cousinweb}

\begin{grammarlens}[title={What you've built}]
One more word for everyday life at A2.
\end{grammarlens}

\subsection*{Your turn}
\begin{itemize}
\raggedright
  \item \emph{Say it:} \emph{tōchaku}
  \item \emph{Say it:} \emph{tōchaku}, once more
  \item \emph{Say it:} \emph{tōchaku wa yūgata desu} — Arrival is in the evening.
  \item \emph{Recall:} say the Japanese for to get to, then the Japanese for departure, then say \emph{tōchaku} again
\end{itemize}

\begin{tcolorbox}[breakable,colback=teal!4,colframe=teal!35!black,title={Before you move on}]
What does \ja{とうちゃく} mean? (\textquotedblleft{}Arrival\textquotedblright{}.) Say it once more.
\end{tcolorbox}
\section[\texorpdfstring{\ja{とっきゅう} (tokkyū)}{とっきゅう (tokkyū)}]{\ja{とっきゅう} (tokkyū) — a limited express}

\label{lesson:JA-C199-tokkyuu}



\begin{quote}
Before the new one: say the Japanese for departure, then the Japanese for arrival.
\end{quote}

\subsection*{You'll want to know: \ja{とっきゅう}}
\textbf{\ja{とっきゅう}} — \emph{tokkyū} — \textquotedblleft{}a limited express\textquotedblright{}.

\begin{quote}
\textbf{\ja{とっきゅうの} \ja{しゅっぱつは} \ja{ごごです。}} — \emph{tokkyū no shuppatsu wa gogo desu} — \textquotedblleft{}The limited express leaves in the afternoon.\textquotedblright{}
\end{quote}

\begin{cousinweb}
\textbf{\ja{とっきゅう}} (\emph{tokkyū}) is Sino-Japanese (kango): two Chinese characters read with their on'yomi, readings borrowed from Chinese. They are \emph{toku}, \textquotedblleft{}special\textquotedblright{} (Middle Chinese \emph{dok}) and \emph{kyū}, \textquotedblleft{}urgent\textquotedblright{} (\emph{kip}) (KANJIDIC2; Middle Chinese in Baxter's transcription). The word is a shortening (JMdict), and the two join as \emph{tok-kyū}.
\end{cousinweb}

\begin{grammarlens}[title={What you've built}]
One more word for everyday life at A2.
\end{grammarlens}

\subsection*{Your turn}
\begin{itemize}
\raggedright
  \item \emph{Say it:} \emph{tokkyū}
  \item \emph{Say it:} \emph{tokkyū}, once more
  \item \emph{Say it:} \emph{tokkyū no shuppatsu wa gogo desu} — The limited express leaves in the afternoon.
  \item \emph{Recall:} say the Japanese for departure, then the Japanese for arrival, then say \emph{tokkyū} again
\end{itemize}

\begin{tcolorbox}[breakable,colback=teal!4,colframe=teal!35!black,title={Before you move on}]
What does \ja{とっきゅう} mean? (\textquotedblleft{}A limited express\textquotedblright{}.) Say it once more.
\end{tcolorbox}
\section[\texorpdfstring{\ja{とまる} (tomaru)}{とまる (tomaru)}]{\ja{とまる} (tomaru) — to come to a stop; to stop moving}

\label{lesson:JA-C199-tomaru}



\begin{quote}
Before the new one: say the Japanese for arrival, then the Japanese for a limited express.
\end{quote}

\subsection*{You'll want to know: \ja{とまる}}
\textbf{\ja{とまる}} — \emph{tomaru} — \textquotedblleft{}to come to a stop; to stop moving\textquotedblright{}.

\begin{quote}
\textbf{\ja{でんしゃが} \ja{えきに} \ja{とまります。}} — \emph{densha ga eki ni tomarimasu} — \textquotedblleft{}The train stops at the station.\textquotedblright{}
\end{quote}

Polite \textbf{\ja{とまります}}, \emph{tomarimasu}; a godan verb.

\begin{cousinweb}
\textbf{\ja{とまる}} (\emph{tomaru}) is a native Japanese word (wago), inherited rather than borrowed from Chinese; its kanji was attached later as a kun'yomi, and its older origin is not settled (JMdict, KANJIDIC2).
\end{cousinweb}

\begin{grammarlens}[title={What you've built}]
One more word for everyday life at A2.
\end{grammarlens}

\subsection*{Your turn}
\begin{itemize}
\raggedright
  \item \emph{Say it:} \emph{tomaru}
  \item \emph{Say it:} \emph{tomaru}, once more
  \item \emph{Say it:} \emph{densha ga eki ni tomarimasu} — The train stops at the station.
  \item \emph{Recall:} say the Japanese for arrival, then the Japanese for a limited express, then say \emph{tomaru} again
\end{itemize}

\begin{tcolorbox}[breakable,colback=teal!4,colframe=teal!35!black,title={Before you move on}]
What does \ja{とまる} mean? (\textquotedblleft{}To come to a stop; to stop moving\textquotedblright{}.) Say it once more.
\end{tcolorbox}
\section[\texorpdfstring{\ja{きゅうこう} (kyūkō)}{きゅうこう (kyūkō)}]{\ja{きゅうこう} (kyūkō) — an express (train)}

\label{lesson:JA-C199-kyuukou}



\begin{quote}
Before the new one: say the Japanese for a limited express, then the Japanese for to come to a stop.
\end{quote}

\subsection*{You'll want to know: \ja{きゅうこう}}
\textbf{\ja{きゅうこう}} — \emph{kyūkō} — \textquotedblleft{}an express (train)\textquotedblright{}.

\begin{quote}
\textbf{\ja{きゅうこうは} \ja{この} \ja{えきに} \ja{とまりません。}} — \emph{kyūkō wa kono eki ni tomarimasen} — \textquotedblleft{}The express does not stop at this station.\textquotedblright{}
\end{quote}

\begin{cousinweb}
\textbf{\ja{きゅうこう}} (\emph{kyūkō}) is Sino-Japanese (kango): two Chinese characters read with their on'yomi, readings borrowed from Chinese. They are \emph{kyū}, \textquotedblleft{}urgent\textquotedblright{} (Middle Chinese \emph{kip}) and \emph{kō}, \textquotedblleft{}to go\textquotedblright{} (\emph{haeng}) (KANJIDIC2; Middle Chinese in Baxter's transcription).
\end{cousinweb}

\begin{grammarlens}[title={What you've built}]
One more word for everyday life at A2.
\end{grammarlens}

\subsection*{Your turn}
\begin{itemize}
\raggedright
  \item \emph{Say it:} \emph{kyūkō}
  \item \emph{Say it:} \emph{kyūkō}, once more
  \item \emph{Say it:} \emph{kyūkō wa kono eki ni tomarimasen} — The express does not stop at this station.
  \item \emph{Recall:} say the Japanese for a limited express, then the Japanese for to come to a stop, then say \emph{kyūkō} again
\end{itemize}

\begin{tcolorbox}[breakable,colback=teal!4,colframe=teal!35!black,title={Before you move on}]
What does \ja{きゅうこう} mean? (\textquotedblleft{}An express (train)\textquotedblright{}.) Say it once more.
\end{tcolorbox}
